% Distilled blueprint for Evans, Chapter 5 (Sobolev Spaces).
% Formal declarations, metadata, and citation identifiers are preserved; exposition and exercises are omitted.

Sobolev spaces provide Banach and Hilbert settings for weak derivatives, approximation, embeddings, compactness, and evolution equations.

\section{Holder Spaces}
\label{sec:holder-spaces}

\begin{definition}[Holder seminorm and sup norm]
\label{def:holder-seminorm-sup-norm}
\dcref{ch5:5.1-def-1}
\group{Hölder Spaces}
(i) If $u : U \to \R$ is bounded and continuous, we write
$$\|u\|_{C(U)} := \sup_{x \in U} |u(x)|.$$
(ii) The $\gamma^{\text{th}}$-Holder seminorm of $u : U \to \R$ is
$$[u]_{C^{0,\gamma}(U)} := \sup_{\substack{x,y \in U \\ x \ne y}} \left\{ \frac{|u(x) - u(y)|}{|x - y|^\gamma} \right\},$$
and the $\gamma^{\text{th}}$-Holder norm is
$$\|u\|_{C^{0,\gamma}(\bar U)} := \|u\|_{C(\bar U)} + [u]_{C^{0,\gamma}(\bar U)}.$$
\end{definition}

\begin{definition}[Holder space $C^{k,\gamma}(\bar U)$]
\label{def:holder-space}
\uses{def:holder-seminorm-sup-norm}
\dcref{ch5:5.1-def-2}
\group{Hölder Spaces}

The Holder space
$$C^{k,\gamma}(\bar U)$$
consists of all functions $u \in C^k(\bar U)$ for which the norm
$$(3) \quad \|u\|_{C^{k,\gamma}(\bar U)} := \sum_{|\alpha| \le k} \|D^\alpha u\|_{C(\bar U)} + \sum_{|\alpha| = k} [D^\alpha u]_{C^{0,\gamma}(\bar U)}$$
is finite.
\end{definition}

\begin{theorem}[Holder spaces as function spaces]
\label{thm:holder-space-banach}
\uses{def:holder-space}
\dcref{ch5:5.1-thm-1}
\group{Hölder Spaces}

The space of functions $C^{k,\gamma}(\bar U)$ is a Banach space.
\end{theorem}

\section{Sobolev Spaces}
\label{sec:sobolev-spaces}

\subsection{Weak derivatives}

\begin{definition}[Weak partial derivative]
\label{def:weak-partial-derivative}
\dcref{ch5:5.2-def-1}
\group{Sobolev Spaces}
Suppose $u, v \in L^1_{\text{loc}}(U)$, and $\alpha$ is a multiindex. We say that $v$ is the $\alpha^{\text{th}}$-weak partial derivative of $u$, written
$$D^\alpha u = v,$$
provided
$$(3) \quad \int_U u D^\alpha \phi \, dx = (-1)^{|\alpha|} \int_U v\phi \, dx$$
for all test functions $\phi \in C_c^\infty(U)$.
\end{definition}

\begin{lemma}[Uniqueness of weak derivatives]
\label{lem:weak-derivative-uniqueness}
\uses{def:weak-partial-derivative}
\dcref{ch5:5.2-lem-1}
\group{Sobolev Spaces}

A weak $\alpha^{\text{th}}$-partial derivative of $u$, if it exists, is uniquely defined up to a set of measure zero.
\end{lemma}

\begin{proof}
\sketch Subtract the two weak identities; the resulting locally integrable function pairs to zero with every test function, hence vanishes almost everywhere.
\end{proof}

\begin{example}[A weak derivative that exists]
\label{ex:weak-derivative-piecewise-linear-example}
\uses{def:weak-partial-derivative}
\dcref{ch5:5.2-ex-1}
\group{Sobolev Spaces}

Let $n = 1$, $U = (0,2)$, and
$$u(x) = \begin{cases}
x & \text{if } 0 < x \le 1 \\
1 & \text{if } 1 \le x < 2.
\end{cases}$$
Define
$$v(x) = \begin{cases}
1 & \text{if } 0 < x \le 1 \\
0 & \text{if } 1 < x < 2.
\end{cases}$$
Let us show $u' = v$ in the weak sense. To see this, choose any $\phi \in C_c^\infty(U)$. We must demonstrate
$$\int_0^2 u\phi' \, dx = -\int_0^2 v\phi \, dx.$$
But we easily calculate
$$\int_0^2 u\phi' \, dx = \int_0^1 x\phi' \, dx + \int_1^2 \phi' \, dx$$
$$= -\int_0^1 \phi \, dx + \phi(1) - \phi(1) = -\int_0^2 v\phi \, dx,$$
as required.
\end{example}

\begin{example}[A weak derivative that fails to exist]
\label{ex:weak-derivative-nonexistence-example}
\uses{def:weak-partial-derivative}
\dcref{ch5:5.2-ex-2}
\group{Sobolev Spaces}

Let $n = 1$, $U = (0,2)$, and
$$u(x) = \begin{cases} x & \text{if } 0 < x \le 1 \\ 2 & \text{if } 1 < x < 2. \end{cases}$$
We assert $u'$ does not exist in the weak sense. To check this, we must show there does not exist any function $v \in L^1_{\text{loc}}(U)$ satisfying
$$(5) \quad \int_0^2 u\phi' \, dx = -\int_0^2 v\phi \, dx$$
for all $\phi \in C_c^\infty(U)$. Suppose, to the contrary, (5) were valid for some $v$ and all $\phi$. Then
$$(6) \quad -\int_0^2 v\phi \, dx = \int_0^2 u\phi' \, dx = \int_0^1 x\phi' \, dx + 2\int_1^2 \phi' \, dx = -\int_0^1 \phi \, dx - \phi(1).$$
Choose a sequence $\{\phi_m\}_{m=1}^\infty$ of smooth functions satisfying
$$0 \le \phi_m \le 1, \quad \phi_m(1) = 1, \quad \phi_m(x) \to 0 \text{ for all } x \ne 1.$$
Replacing $\phi$ by $\phi_m$ in (6) and sending $m \to \infty$, we discover
$$1 = \lim_{m \to \infty} \phi_m(1) = \lim_{m \to \infty} \left[\int_0^2 v\phi_m \, dx - \int_0^1 \phi_m \, dx\right] = 0,$$
a contradiction.
\end{example}

\subsection{Definition of Sobolev spaces}

\begin{definition}[Sobolev space $W^{k,p}(U)$]
\label{def:sobolev-space-wkp}
\uses{def:weak-partial-derivative}
\dcref{ch5:5.2-def-2}
\group{Sobolev Spaces}

The Sobolev space
$$W^{k,p}(U)$$
consists of all locally summable functions $u : U \to \R$ such that for each multiindex $\alpha$ with $|\alpha| \le k$, $D^\alpha u$ exists in the weak sense and belongs to $L^p(U)$.
\end{definition}

\begin{remark}[Notation $H^k$ and identification a.e.]
\label{rem:hk-notation-identification}
\uses{def:sobolev-space-wkp}
\dcref{ch5:5.2-rem-1}
\group{Sobolev Spaces}

(i) If $p = 2$, we usually write
$$H^k(U) = W^{k,2}(U) \quad (k = 0, 1, \ldots).$$
In particular, $H^0(U) = L^2(U)$.

(ii) Functions in $W^{k,p}(U)$ that agree a.e. are identified.
\end{remark}

\begin{definition}[Norm on $W^{k,p}(U)$]
\label{def:sobolev-norm}
\uses{def:sobolev-space-wkp}
\dcref{ch5:5.2-def-3}
\group{Sobolev Spaces}

If $u \in W^{k,p}(U)$, we define its norm to be
$$\|u\|_{W^{k,p}(U)} := \begin{cases}
\left(\sum_{|\alpha| \le k} \int_U |D^\alpha u|^p\,dx\right)^{1/p} & (1 \le p < \infty) \\
\sum_{|\alpha| \le k} \text{ess sup}_U |D^\alpha u| & (p = \infty).
\end{cases}$$
\end{definition}

\begin{definition}[Convergence in $W^{k,p}(U)$ and $W^{k,p}_{\text{loc}}(U)$]
\label{def:sobolev-convergence}
\uses{def:sobolev-norm}
\dcref{ch5:5.2-def-4}
\group{Sobolev Spaces}

(i) Let $\{u_m\}_{m=1}^\infty$, $u \in W^{k,p}(U)$. We say $u_m$ converges to $u$ in $W^{k,p}(U)$, written
$$u_m \to u \quad \text{in } W^{k,p}(U),$$
provided
$$\lim_{m \to \infty} \|u_m - u\|_{W^{k,p}(U)} = 0.$$

(ii) We write
$$u_m \to u \quad \text{in } W^{k,p}_{\text{loc}}(U),$$
to mean
$$u_m \to u \quad \text{in } W^{k,p}(V)$$
for each $V \CC U$.
\end{definition}

\begin{definition}[The space $W_0^{k,p}(U)$]
\label{def:sobolev-space-zero-boundary}
\uses{def:sobolev-convergence}
\dcref{ch5:5.2-def-5}
\group{Sobolev Spaces}

We denote by
$$W_0^{k,p}(U)$$
the closure of $C_c^\infty(U)$ in $W^{k,p}(U)$.
\end{definition}

\begin{example}[A power singularity in a Sobolev space]
\label{ex:power-singularity-sobolev-membership}
\uses{def:sobolev-space-wkp, def:weak-partial-derivative}
\dcref{ch5:5.2-ex-3}
\group{Sobolev Spaces}

Take $U = B^0(0,1)$, the open unit ball in $\R^n$, and
$$u(x) = |x|^{-\alpha} \quad (x \in U, x \ne 0).$$
For which values of $\alpha > 0, n, p$ does $u$ belong to $W^{1,p}(U)$? To answer, note first $u$ is smooth away from $0$, with
$$u_{x_i}(x) = \frac{-\alpha x_i}{|x|^{\alpha+2}} \quad (x \ne 0),$$
and so
$$|Du(x)| = \frac{\alpha}{|x|^{\alpha+1}} \quad (x \ne 0).$$
Let $\phi \in C_c^\infty(U)$ and fix $\varepsilon > 0$. Then
$$\int_{U - B(0,\varepsilon)} u\phi_{x_i}\,dx = -\int_{U - B(0,\varepsilon)} u_{x_i}\phi\,dx + \int_{\partial B(0,\varepsilon)} u\phi\nu^i\,dS,$$
$\nu = (\nu^1, \ldots, \nu^n)$ denoting the inward pointing normal on $\partial B(0,\varepsilon)$. Now if $\alpha + 1 < n$, $|Du(x)| \in L^1(U)$. In this case
$$\left|\int_{\partial B(0,\varepsilon)} u\phi\nu^i\,dS\right| \le \|\phi\|_{L^\infty} \int_{\partial B(0,\varepsilon)} \varepsilon^{-\alpha}\,dS \le C\varepsilon^{n-1-\alpha} \to 0.$$
Thus
$$\int_U u\phi_{x_i}\,dx = -\int_U u_{x_i}\phi\,dx$$
for all $\phi \in C_c^\infty(U)$, provided $0 \le \alpha < n - 1$. Furthermore $|Du(x)| = \frac{\alpha}{|x|^{\alpha+1}} \in L^p(U)$ if and only if $(\alpha+1)p < n$. Consequently $u \in W^{1,p}(U)$ if and only if $\alpha < \frac{n-p}{p}$. In particular $u \notin W^{1,p}(U)$ for each $p \ge n$.
\end{example}

\begin{example}[A dense countable sum of singularities]
\label{ex:dense-singularities-unbounded-sobolev}
\uses{ex:power-singularity-sobolev-membership}
\dcref{ch5:5.2-ex-4}
\group{Sobolev Spaces}

Let $\{r_k\}_{k=1}^{\infty}$ be a countable, dense subset of $U = B^0(0,1)$. Write
$$u(x) = \sum_{k=1}^{\infty} \frac{1}{2^k} |x - r_k|^{-\alpha} \quad (x \in U).$$
Then $u \in W^{1,p}(U)$ if and only if $\alpha < \frac{n-p}{p}$. If $0 < \alpha < \frac{n-p}{p}$, we see that $u$ belongs to $W^{1,p}(U)$ and yet is unbounded on each open subset of $U$.
\end{example}

\subsection{Elementary properties}

\begin{theorem}[Properties of weak derivatives]
\label{thm:weak-derivative-properties}
\uses{def:sobolev-space-wkp}
\dcref{ch5:5.2-thm-1}
\group{Sobolev Spaces}

Assume $u, v \in W^{k,p}(U)$, $|\alpha| \le k$. Then

(i) $D^{\alpha}u \in W^{k-|\alpha|,p}(U)$ and $D^{\beta}(D^{\alpha}u) = D^{\alpha}(D^{\beta}u) = D^{\alpha+\beta}u$ for all multiindices $\alpha, \beta$ with $|\alpha| + |\beta| \le k$.

(ii) For each $\lambda, \mu \in \R$, $\lambda u + \mu v \in W^{k,p}(U)$ and $D^{\alpha}(\lambda u + \mu v) = \lambda D^{\alpha}u + \mu D^{\alpha}v$, $|\alpha| \le k$.

(iii) If $V$ is an open subset of $U$, then $u \in W^{k,p}(V)$.

(iv) If $\zeta \in C_c^{\infty}(U)$, then $\zeta u \in W^{k,p}(U)$ and
$$(7) \quad D^{\alpha}(\zeta u) = \sum_{\beta \le \alpha} \binom{\alpha}{\beta} D^{\beta}\zeta\, D^{\alpha-\beta}u \quad (\textit{Leibniz' formula}),$$
where $\binom{\alpha}{\beta} = \frac{\alpha!}{\beta!(\alpha-\beta)!}$.
\end{theorem}

\begin{proof}
\sketch Test the iterated weak derivative identity against derivatives of a test function. Linearity and restriction are immediate, while induction on $|\alpha|$ gives the Leibniz formula.
\end{proof}

\begin{theorem}[Sobolev spaces as function spaces]
\label{thm:sobolev-space-banach}
\uses{def:sobolev-norm}
\dcref{ch5:5.2-thm-2}
\group{Sobolev Spaces}

For each $k = 1, \ldots$ and $1 \le p \le \infty$, the Sobolev space $W^{k,p}(U)$ is a Banach space.
\end{theorem}

\begin{proof}
\sketch The norm properties follow from the $L^p$ triangle inequality. A Cauchy sequence has $L^p$ limits for every weak derivative; passing to the defining distributional identities identifies these limits as the derivatives of the $L^p$ limit.
\end{proof}

\section{Approximation}
\label{sec:sobolev-approximation}

\subsection{Interior approximation by smooth functions}

\begin{theorem}[Local approximation by smooth functions]
\label{thm:local-smooth-approximation}
\uses{def:sobolev-space-wkp, def:sobolev-convergence}
\dcref{ch5:5.3-thm-1}
\group{Sobolev Space Approximation}

Assume $u \in W^{k,p}(U)$ for some $1 \le p < \infty$, and set
$$u^\varepsilon = \eta_\varepsilon * u \text{ in } U_\varepsilon.$$
Then
\begin{enumerate}
\item[(i)] $u^\varepsilon \in C^\infty(U_\varepsilon)$ for each $\varepsilon > 0$,
\end{enumerate}
and
\begin{enumerate}
\item[(ii)] $u^\varepsilon \to u$ in $W^{k,p}_{\text{loc}}(U)$, as $\varepsilon \to 0$.
\end{enumerate}
\end{theorem}

\begin{proof}
\sketch Mollification is smooth on $U_\varepsilon$ and satisfies $D^\alpha u^\varepsilon=\eta_\varepsilon*D^\alpha u$. Approximate-identity convergence in $L^p$ for each $|\alpha|\le k$ gives local Sobolev convergence.
\end{proof}

\subsection{Global approximation by smooth functions}

\begin{theorem}[Global approximation by smooth functions]
\label{thm:global-smooth-approximation-interior}
\uses{def:sobolev-space-wkp, thm:weak-derivative-properties}
\dcref{ch5:5.3-thm-2}
\group{Sobolev Space Approximation}

Assume $U$ is bounded, and suppose as well that $u \in W^{k,p}(U)$ for some $1 \le p < \infty$. Then there exist functions $u_m \in C^\infty(U) \cap W^{k,p}(U)$ such that
$$u_m \to u \quad \text{in } W^{k,p}(U).$$
\end{theorem}

\begin{proof}
\sketch Use a locally finite smooth partition of unity, mollify each compactly supported piece with errors summing to a prescribed $\delta$, and sum the mollified pieces. Local finiteness gives smoothness and the geometric error bound gives global convergence.
\end{proof}

\begin{remark}[Approximation not up to the boundary]
\label{rem:approximation-not-up-to-boundary}
\uses{thm:global-smooth-approximation-interior}
\dcref{ch5:5.3-rem-1}
\group{Sobolev Space Approximation}

The approximants need not belong to $C^\infty(\bar U)$.
\end{remark}

\subsection{Global approximation by functions smooth up to the boundary}

\begin{theorem}[Global approximation by functions smooth up to the boundary]
\label{thm:global-smooth-approximation-boundary}
\uses{def:sobolev-space-wkp, thm:local-smooth-approximation, thm:weak-derivative-properties}
\dcref{ch5:5.3-thm-3}
\group{Sobolev Space Approximation}

Assume $U$ is bounded and $\partial U$ is $C^1$. Suppose $u \in W^{k,p}(U)$ for some $1 \le p < \infty$. Then there exist functions $u_m \in C^\infty(\bar U)$ such that
$$u_m \to u \quad \text{in } W^{k,p}(U).$$
\end{theorem}

\begin{proof}
\sketch Flatten the $C^1$ boundary, translate inward before mollifying, and use translation plus approximate-identity convergence in each chart. A finite partition of unity patches the chartwise approximants and preserves the Sobolev estimate.
\end{proof}

\section{Extensions}
\label{sec:sobolev-extensions}

\begin{theorem}[Extension Theorem]
\label{thm:sobolev-extension-theorem}
\uses{def:sobolev-space-wkp, thm:global-smooth-approximation-boundary}
\dcref{ch5:5.4-thm-1}
\group{Sobolev Extension}

Assume $U$ is bounded and $\partial U$ is $C^1$. Select a bounded open set $V$ such that $U \CC V$. Then there exists a bounded linear operator
$$(1) \quad E: W^{1,p}(U) \to W^{1,p}(\R^n)$$
such that for each $u \in W^{1,p}(U)$:
\begin{itemize}
\item[(i)] $Eu = u$ a.e. in $U$,
\item[(ii)] $Eu$ has support within $V$,
\end{itemize}
and
\begin{itemize}
\item[(iii)] $\|Eu\|_{W^{1,p}(\R^n)} \le C\|u\|_{W^{1,p}(U)},$
\end{itemize}
the constant $C$ depending only on $p$, $U$, and $V$.
\end{theorem}

\begin{proof}
\sketch Reflect across a flattened boundary by a finite linear combination of inward values chosen to match first derivatives, then multiply by a cutoff. Change of variables and a finite boundary cover give a bounded global operator with support in $V$; density extends the construction from smooth functions to $W^{1,p}$.
\end{proof}

\begin{definition}[Extension of a Sobolev function]
\label{def:extension-of-sobolev-function}
\uses{thm:sobolev-extension-theorem}
\dcref{ch5:5.4-def-1}
\group{Sobolev Extension}

We call $Eu$ an extension of $u$ to $\R^n$.
\end{definition}

\begin{remark}[Extension in $W^{2,p}$; failure for $k>2$]
\label{rem:extension-higher-order-remark}
\uses{thm:sobolev-extension-theorem}
\dcref{ch5:5.4-rem-1}
\group{Sobolev Extension}

(i) If $\partial U$ is $C^2$, the same operator is bounded from $W^{2,p}(U)$ to $W^{2,p}(\R^n)$ and satisfies
$$(10) \quad \|Eu\|_{W^{2,p}(\R^n)} \le C\|u\|_{W^{2,p}(U)},$$
where $C$ depends only on $U$, $V$, $n$, and $p$.

(ii) This reflection construction does not provide extensions in $W^{k,p}(U)$ for $k>2$.
\end{remark}

\section{Traces}
\label{sec:sobolev-traces}

\begin{theorem}[Trace Theorem]
\label{thm:trace-theorem}
\uses{def:sobolev-space-wkp, thm:global-smooth-approximation-boundary}
\dcref{ch5:5.5-thm-1}
\group{Sobolev Traces}

Assume $U$ is bounded and $\partial U$ is $C^1$. Then there exists a bounded linear operator
$$T : W^{1,p}(U) \to L^p(\partial U)$$
such that
\begin{itemize}
\item[(i)] $Tu = u|_{\partial U}$ if $u \in W^{1,p}(U) \cap C(\overline{U})$
\end{itemize}
and
\begin{itemize}
\item[(ii)] $$\|Tu\|_{L^p(\partial U)} \le C\|u\|_{W^{1,p}(U)},$$
\end{itemize}
for each $u \in W^{1,p}(U)$, with the constant $C$ depending only on $p$ and $U$.
\end{theorem}

\begin{proof}
\sketch For smooth functions, integrate the normal derivative across flattened boundary charts and apply Young inequality to obtain the local boundary estimate. Patch finitely many charts, then define the trace by completing smooth approximants in $L^p(\partial U)$.
\end{proof}

\begin{definition}[Trace of $u$ on $\partial U$]
\label{def:trace-operator}
\uses{thm:trace-theorem}
\dcref{ch5:5.5-def-1}
\group{Sobolev Traces}

We call $Tu$ the trace of $u$ on $\partial U$.
\end{definition}

\begin{theorem}[Trace-zero functions in $W^{1,p}$]
\label{thm:trace-zero-characterization}
\uses{def:sobolev-space-zero-boundary, thm:trace-theorem}
\dcref{ch5:5.5-thm-2}
\group{Sobolev Traces}

Assume $U$ is bounded and $\partial U$ is $C^1$. Suppose furthermore that $u \in W^{1,p}(U)$. Then
$$(4) \quad u \in W_0^{1,p}(U) \quad \text{if and only if} \quad Tu = 0 \text{ on } \partial U.$$
\end{theorem}

\begin{proof}
\sketch The forward implication follows from continuity of the trace on the closure of $C_c^\infty(U)$. For the converse, flatten the boundary, use the zero trace to control the boundary layer by $Du$, and cut off that layer; the cutoff error tends to zero in $W^{1,p}$.
\end{proof}

\section{Sobolev Inequalities}
\label{sec:sobolev-inequalities}

\subsection{Gagliardo--Nirenberg--Sobolev inequality}

\begin{definition}[Sobolev conjugate exponent]
\label{def:sobolev-conjugate-exponent}
\dcref{ch5:5.6-def-1}
\group{Sobolev Inequalities}
If $1 \le p < n$, the Sobolev conjugate of $p$ is
$$(8) \quad p^* := \frac{np}{n - p}.$$
Note that
$$(9) \quad \frac{1}{p^*} = \frac{1}{p} - \frac{1}{n}, \quad p^* > p.$$
\end{definition}

\begin{theorem}[Gagliardo--Nirenberg--Sobolev inequality]
\label{thm:gagliardo-nirenberg-sobolev-inequality}
\uses{def:sobolev-conjugate-exponent}
\dcref{ch5:5.6-thm-1}
\group{Sobolev Inequalities}

Assume $1 \le p < n$. There exists a constant $C$, depending only on $p$ and $n$, such that
$$(10) \quad \|u\|_{L^{p^*}(\R^n)} \le C\|Du\|_{L^p(\R^n)},$$
for all $u \in C_c^1(\R^n)$.
\end{theorem}

\begin{proof}
\sketch For compactly supported smooth $u$, represent $u$ by the Newtonian kernel and apply Holder inequality with the integrability of $|x|^{1-n}$. Approximation by $C_c^\infty$ extends the estimate to $W^{1,p}$.
\end{proof}

\begin{theorem}[Estimates for $W^{1,p}$, $1 \le p < n$]
\label{thm:sobolev-embedding-w1p-subcritical}
\uses{thm:gagliardo-nirenberg-sobolev-inequality, thm:sobolev-extension-theorem, thm:local-smooth-approximation}
\dcref{ch5:5.6-thm-2}
\group{Sobolev Inequalities}

Let $U$ be a bounded, open subset of $\R^n$, and suppose $\partial U$ is $C^1$. Assume $1 \le p < n$, and $u \in W^{1,p}(U)$. Then $u \in L^{p^*}(U)$, with the estimate
$$(15) \quad \|u\|_{L^{p^*}(U)} \le C\|u\|_{W^{1,p}(U)},$$
the constant $C$ depending only on $p, n$, and $U$.
\end{theorem}

\begin{proof}
\sketch Extend $u$ to a compactly supported function on $\mathbb R^n$, approximate by smooth functions, and apply the global inequality to the approximants. Restriction and the extension bound yield the stated $L^q$ estimate.
\end{proof}

\begin{theorem}[Poincare's inequality for $W_0^{1,p}$, $1 \le p < n$]
\label{thm:poincare-w01p-subcritical}
\uses{def:sobolev-space-zero-boundary, thm:gagliardo-nirenberg-sobolev-inequality}
\dcref{ch5:5.6-thm-3}
\group{Sobolev Inequalities}

Assume $U$ is a bounded, open subset of $\R^n$. Suppose $u \in W_0^{1,p}(U)$ for some $1 \le p < n$. Then we have the estimate
$$\|u\|_{L^q(U)} \le C\|Du\|_{L^p(U)}$$
for each $q \in [1, p^*]$, the constant $C$ depending only on $p, q, n$ and $U$.
\end{theorem}

\begin{proof}
\sketch Approximate by compactly supported smooth functions, extend them by zero, and apply the global Gagliardo--Nirenberg--Sobolev estimate. The finite measure of $U$ gives all exponents $q\le p^*$.
\end{proof}

\begin{remark}[Poincare's inequality]
\label{rem:poincare-w01p-terminology}
\uses{thm:poincare-w01p-subcritical, thm:sobolev-embedding-w1p-subcritical}
\dcref{ch5:5.6-rem-1}
\group{Sobolev Inequalities}

The estimate in \cref{thm:poincare-w01p-subcritical} is Poincare's inequality; unlike \cref{thm:sobolev-embedding-w1p-subcritical}, its right-hand side contains only $Du$.
\end{remark}

\begin{remark}[Equivalent norm on $W_0^{1,p}$; the borderline case $p=n$]
\label{rem:borderline-case-p-equals-n}
\uses{thm:poincare-w01p-subcritical}
\dcref{ch5:5.6-rem-2}
\group{Sobolev Inequalities}

(i) If $U$ is bounded, $\|Du\|_{L^p(U)}$ and $\|u\|_{W^{1,p}(U)}$ are equivalent norms on $W_0^{1,p}(U)$.

(ii) For $n>1$, the borderline embedding $W^{1,n}(U)\subset L^\infty(U)$ fails. On $U=B^0(0,1)$,
$$u(x)=\log\log\left(1+\frac{1}{|x|}\right)$$
belongs to $W^{1,n}(U)$ but not to $L^\infty(U)$.
\end{remark}

\subsection{Morrey's inequality}

\begin{theorem}[Morrey's inequality]
\label{thm:morreys-inequality}
\uses{def:holder-seminorm-sup-norm}
\dcref{ch5:5.6-thm-4}
\group{Sobolev Inequalities}

Assume $n < p \le \infty$. Then there exists a constant $C$, depending only on $p$ and $n$, such that
$$(20) \quad \|u\|_{C^{0,\gamma}(\R^n)} \le C\|u\|_{W^{1,p}(\R^n)}$$
for all $u \in C^1(\R^n)$, where
$$\gamma := 1 - n/p.$$
\end{theorem}

\begin{proof}
\sketch Radial integration gives an average oscillation bound by the Riesz kernel. Holder inequality yields the uniform bound and the estimate $|u(x)-u(y)|\le C|x-y|^{1-n/p}\|Du\|_{L^p}$.
\end{proof}

\begin{remark}[Morrey estimate on a ball]
\label{rem:morrey-estimate-on-ball}
\uses{thm:morreys-inequality}
\dcref{ch5:5.6-rem-3}
\group{Sobolev Inequalities}

For $n<p<\infty$, $u\in W^{1,p}(B(x,2r))$, and $y\in B(x,r)$,
$$|u(y) - u(x)| \le Cr^{1-n/p} \left( \int_{B(x,2r)} |Du(z)|^p\,dz \right)^{1/p}$$
with $C$ depending only on $n$ and $p$. The integral may also be restricted to $B(x,r)$.
\end{remark}

\begin{definition}[Version of a function]
\label{def:version-of-a-function}
\dcref{ch5:5.6-def-2}
\group{Sobolev Inequalities}
We say $u^*$ is a version of a given function $u$ provided
$$u = u^* \text{ a.e.}$$
\end{definition}

\begin{theorem}[Estimates for $W^{1,p}$, $n < p \le \infty$]
\label{thm:sobolev-embedding-w1p-supercritical}
\uses{def:version-of-a-function, thm:morreys-inequality, thm:sobolev-extension-theorem, thm:local-smooth-approximation}
\dcref{ch5:5.6-thm-5}
\group{Sobolev Inequalities}

Let $U$ be a bounded, open subset of $\R^n$, and suppose $\partial U$ is $C^1$. Assume $n < p \le \infty$, and $u \in W^{1,p}(U)$. Then $u$ has a version $u^* \in C^{0,\gamma}(\bar{U})$, for $\gamma = 1 - \frac{n}{p}$, with the estimate
$$\|u^*\|_{C^{0,\gamma}(\bar{U})} \le C\|u\|_{W^{1,p}(U)}.$$
The constant $C$ depends only on $p, n$ and $U$.
\end{theorem}

\begin{proof}
\sketch Extend to compact support in $\mathbb R^n$, mollify, and use Morrey estimate to obtain a Cauchy sequence in the Holder norm. Its limit is the continuous version and inherits the extension bound.
\end{proof}

\begin{remark}[Identification with continuous version]
\label{rem:identify-continuous-version}
\uses{thm:sobolev-embedding-w1p-supercritical}
\dcref{ch5:5.6-rem-4}
\group{Sobolev Inequalities}

For $p>n$, each $u\in W^{1,p}(U)$ is identified with the continuous version supplied by \cref{thm:sobolev-embedding-w1p-supercritical}.
\end{remark}

\subsection{General Sobolev inequalities}

\begin{theorem}[General Sobolev inequalities]
\label{thm:general-sobolev-inequalities}
\uses{thm:gagliardo-nirenberg-sobolev-inequality, thm:morreys-inequality}
\dcref{ch5:5.6-thm-6}
\group{Sobolev Inequalities}

Let $U$ be a bounded open subset of $\R^n$, with a $C^1$ boundary. Assume $u \in W^{k,p}(U)$.

(i) If
$$(29) \quad k < \frac{n}{p},$$
then $u \in L^q(U)$, where
$$\frac{1}{q} = \frac{1}{p} - \frac{k}{n}.$$
We have in addition the estimate
$$(30) \quad \|u\|_{L^q(U)} \le C\|u\|_{W^{k,p}(U)},$$
the constant $C$ depending only on $k, p, n$ and $U$.

(ii) If
$$(31) \quad k > \frac{n}{p},$$
then $u \in C^{k-[\frac{n}{p}]-1,\gamma}(\bar{U})$, where
$$\gamma = \begin{cases}
\left[\frac{n}{p}\right] + 1 - \frac{n}{p}, & \text{if } \frac{n}{p} \text{ is not an integer} \\
\text{any positive number} < 1, & \text{if } \frac{n}{p} \text{ is an integer}.
\end{cases}$$
We have in addition the estimate
$$(32) \quad \|u\|_{C^{k-[\frac{n}{p}]-1,\gamma}(\bar{U})} \le C\|u\|_{W^{k,p}(U)},$$
the constant $C$ depending only on $k, p, n, \gamma$ and $U$.
\end{theorem}

\begin{proof}
\sketch Iterate the subcritical Sobolev embedding until reaching $L^q$ when $k<n/p$. When $k>n/p$, iterate to an exponent above $n$ and apply Morrey; at the integer borderline use every finite exponent $q>n$.
\end{proof}

\begin{remark}[Fourier-transform proof]
\label{rem:general-sobolev-fourier-remark}
\uses{thm:general-sobolev-inequalities}
\dcref{ch5:5.6-rem-5}
\group{Sobolev Inequalities}

General Sobolev inequalities also admit Fourier-transform proofs.
\end{remark}

\section{Compactness}
\label{sec:sobolev-compactness}

\begin{definition}[Compact embedding]
\label{def:compact-embedding}
\dcref{ch5:5.7-def-1}
\group{Compactness}
Let $X$ and $Y$ be Banach spaces, $X \subset Y$. We say that $X$ is compactly embedded in $Y$, written
$$X \CC Y,$$
provided
\begin{itemize}
\item[(i)] $\|x\|_Y \le C\|x\|_X$ ($x \in X$) for some constant $C$,
\end{itemize}
and

(ii) each bounded sequence in $X$ is precompact in $Y$.
\end{definition}

\begin{theorem}[Rellich--Kondrachov Compactness Theorem]
\label{thm:rellich-kondrachov}
\uses{def:compact-embedding, thm:sobolev-embedding-w1p-subcritical, thm:sobolev-extension-theorem, thm:local-smooth-approximation}
\dcref{ch5:5.7-thm-1}
\group{Compactness}

Assume $U$ is a bounded open subset of $\R^n$, and $\partial U$ is $C^1$. Suppose $1 \le p < n$. Then
$$W^{1,p}(U) \CC L^q(U)$$
for each $1 \le q < p^*$.
\end{theorem}

\begin{proof}
\sketch After extension, assume compact support in a fixed bounded set. Mollification converges uniformly in $L^q$ on bounded Sobolev sets, while for fixed mollifier scale the sequence is bounded and equicontinuous; Arzela--Ascoli followed by a diagonal argument gives a convergent subsequence.
\end{proof}

\begin{remark}[$W^{1,p}(U) \CC L^p(U)$ for all $p$]
\label{rem:compact-embedding-lp-all-p}
\uses{thm:rellich-kondrachov, thm:sobolev-embedding-w1p-subcritical, thm:morreys-inequality}
\dcref{ch5:5.7-rem-1}
\group{Compactness}

For every $1\le p\le\infty$,
$$W^{1,p}(U) \CC L^p(U)$$
when $U$ is bounded with $C^1$ boundary. Moreover,
$$W_0^{1,p}(U) \CC L^p(U),$$
without a $C^1$ boundary assumption.
\end{remark}

\section{Additional Topics}
\label{sec:sobolev-additional-topics}

\subsection{Poincare's inequalities}

\begin{theorem}[Poincare's inequality]
\label{thm:poincares-inequality-connected}
\uses{thm:rellich-kondrachov, def:sobolev-space-wkp}
\dcref{ch5:5.8-thm-1}
\group{Poincaré Inequalities}

Let $U$ be a bounded, connected, open subset of $\R^n$, with a $C^1$ boundary $\partial U$. Assume $1 \le p \le \infty$. Then there exists a constant $C$, depending only on $n,p$ and $U$, such that
$$(1) \quad \|u - (u)_U\|_{L^p(U)} \le C\|Du\|_{L^p(U)}$$
for each function $u \in W^{1,p}(U)$.
\end{theorem}

\begin{proof}
\sketch If the estimate failed, normalize a sequence with vanishing gradients. Rellich compactness gives an $L^p$ limit, the limit has zero weak gradient and is constant on connected $U$, while the normalization contradicts convergence.
\end{proof}

\begin{theorem}[Poincare's inequality for a ball]
\label{thm:poincares-inequality-ball}
\uses{thm:poincares-inequality-connected}
\dcref{ch5:5.8-thm-2}
\group{Poincaré Inequalities}

Assume $1 \le p \le \infty$. Then there exists a constant $C$, depending only on $n$ and $p$, such that
$$(7) \quad \|u - (u)_{x,r}\|_{L^p(B(x,r))} \le Cr\|Du\|_{L^p(B(x,r))}$$
for each ball $B(x,r) \subset \R^n$ and each function $u \in W^{1,p}(B^0(x,r))$.
\end{theorem}

\begin{proof}
\sketch Apply the connected-domain inequality to a ball and subtract the mean. Scaling gives the stated radius dependence and the $L^p$ estimate.
\end{proof}

\begin{remark}[Bounded mean oscillation]
\label{rem:bmo-remark}
\uses{thm:poincares-inequality-ball}
\dcref{ch5:5.8-rem-1}
\group{Poincaré Inequalities}

If $u \in W^{1,n}(\R^n) \cap L^1(\R^n)$ and $B(x,r)$ is a ball, then
$$\int_{B(x,r)} |u - (u)_{x,r}|\, dy \le Cr\int_{B(x,r)} |Du|\, dy$$
$$\le Cr \left(\int_{B(x,r)} |Du|^n\, dy\right)^{1/n} \le C\left(\int_{\R^n} |Du|^n\, dy\right)^{1/n}.$$
Thus $u \in BMO(\R^n)$, with seminorm
$$[u]_{BMO(\R^n)} := \sup_{B(x,r) \subset \R^n} \left\{\int_{B(x,r)} |u - (u)_{x,r}|\, dy\right\}.$$
See Stein \cite{Stein1993} for the theory of $BMO$.
\end{remark}

\subsection{Difference quotients}

\begin{definition}[Difference quotients]
\label{def:difference-quotients}
\dcref{ch5:5.8-def-1}
\group{Difference Quotients}
(i) The $i^{\text{th}}$-difference quotient of size $h$ is
$$D_i^h u(x) = \frac{u(x + he_i) - u(x)}{h} \quad (i = 1, \ldots, n)$$
for $x \in V$ and $h \in \R$, $0 < |h| < \dist(V, \partial U)$.

(ii) $D^h u := (D_1^h u, \ldots, D_n^h u)$.
\end{definition}

\begin{theorem}[Difference quotients and weak derivatives]
\label{thm:difference-quotients-weak-derivatives}
\uses{def:difference-quotients, def:sobolev-space-wkp}
\dcref{ch5:5.8-thm-3}
\group{Difference Quotients}

(i) Suppose $1 \le p < \infty$ and $u \in W^{1,p}(U)$. Then for each $V \CC U$
$$(8) \quad \|D^h u\|_{L^p(V)} \le C\|Du\|_{L^p(U)}$$
for some constant $C$ and all $0 < |h| < \frac{1}{2}\dist(V, \partial U)$.

(ii) Assume $1 < p < \infty$, $u \in L^p(V)$, and there exists a constant $C$ such that
$$(9) \quad \|D^h u\|_{L^p(V)} \le C$$
for all $0 < |h| < \frac{1}{2}\dist(V, \partial U)$. Then
$$u \in W^{1,p}(V), \quad \text{with } \|Du\|_{L^p(V)} \le C.$$
\end{theorem}

\begin{proof}
\sketch The difference-quotient identity follows by translating test functions. Uniform $L^p$ bounds yield weakly convergent difference quotients; passing to the limit in the identity identifies the limit with the weak derivative.
\end{proof}

\begin{remark}[Assertion (ii) fails for $p=1$]
\label{rem:difference-quotient-p1-fails}
\uses{thm:difference-quotients-weak-derivatives}
\dcref{ch5:5.8-rem-2}
\group{Difference Quotients}

Assertion (ii) of \cref{thm:difference-quotients-weak-derivatives} is false for $p=1$.
\end{remark}

\begin{remark}[Tangential difference quotients on half-balls]
\label{rem:difference-quotient-half-ball-remark}
\uses{thm:difference-quotients-weak-derivatives}
\dcref{ch5:5.8-rem-3}
\group{Difference Quotients}

For the half-balls
$$U=B^0(0,1)\cap\{x_n>0\},\qquad V=B^0(0,1/2)\cap\{x_n>0\},$$
the tangential difference quotients satisfy
$$\int_V |D_i^h u|^p\,dx \le \int_U |u_{x_i}|^p\,dx \qquad (i=1,\ldots,n-1),$$
although $V\not\CC U$.
\end{remark}

\begin{theorem}[Characterization of $W^{1,\infty}$]
\label{thm:characterization-w1infty-lipschitz}
\uses{def:sobolev-space-wkp, thm:sobolev-extension-theorem}
\dcref{ch5:5.8-thm-4}
\group{Difference Quotients}

Let $U$ be open and bounded, with $\partial U$ of class $C^1$. Then $u : U \to \R$ is Lipschitz continuous if and only if $u \in W^{1,\infty}(U)$.
\end{theorem}

\begin{proof}
\sketch For smooth functions, integrate the gradient along line segments to obtain Lipschitz bounds. Conversely, bounded difference quotients have weakly convergent subsequences, and the translation identity identifies their limits with the weak derivatives.
\end{proof}

\begin{remark}[Local Lipschitz characterization; failure for $1\le p<\infty$]
\label{rem:local-lipschitz-w1infty-loc-remark}
\uses{thm:characterization-w1infty-lipschitz, thm:sobolev-embedding-w1p-supercritical}
\dcref{ch5:5.8-rem-4}
\group{Difference Quotients}

For every open $U$, $u\in W^{1,\infty}_{\mathrm{loc}}(U)$ if and only if $u$ is locally Lipschitz.
No analogous equivalence holds for finite $p$: if $n<p<\infty$, then $W^{1,p}\subset C^{0,1-n/p}$, but Holder continuity with exponent less than one does not imply membership in any $W^{1,p}$.
\end{remark}

\subsection{Differentiability a.e.}

\begin{definition}[Differentiability at a point]
\label{def:differentiability-at-a-point}
\dcref{ch5:5.8-def-2}
\group{Differentiability of Sobolev Functions}
A function $u: U \to \R$ is differentiable at $x \in U$ if there exists $a \in \R^n$ such that
$$(13) \quad u(y) = u(x) + a \cdot (y - x) + o(|y - x|) \quad \text{as } y \to x.$$
In other words,
$$\lim_{y \to x} \frac{|u(y) - u(x) - a \cdot (y - x)|}{|y - x|} = 0.$$
\end{definition}

\begin{theorem}[Differentiability almost everywhere]
\label{thm:differentiability-almost-everywhere}
\uses{def:differentiability-at-a-point, thm:morreys-inequality}
\dcref{ch5:5.8-thm-5}
\group{Differentiability of Sobolev Functions}

Assume $u \in W^{1,p}_{\text{loc}}(U)$ for some $n < p \le \infty$. Then $u$ is differentiable a.e. in $U$, and its gradient equals its weak gradient a.e.
\end{theorem}

\begin{proof}
\sketch Apply the local Morrey estimate to $u(y)-u(x)-Du(x)\cdot(y-x)$ at Lebesgue points of $Du$. The resulting oscillation is $o(|x-y|)$; the case $p=\infty$ follows by inclusion in finite exponents.
\end{proof}

\begin{theorem}[Rademacher's Theorem]
\label{thm:rademachers-theorem}
\uses{thm:differentiability-almost-everywhere, thm:characterization-w1infty-lipschitz}
\dcref{ch5:5.8-thm-6}
\group{Differentiability of Sobolev Functions}

Let $u$ be locally Lipschitz continuous in $U$. Then $u$ is differentiable almost everywhere in $U$.
\end{theorem}

\subsection{Fourier transform methods}

\begin{theorem}[Characterization of $H^k$ by Fourier transform]
\label{thm:characterization-hk-fourier-transform}
\uses{rem:hk-notation-identification}
\dcref{ch5:5.8-thm-7}
\group{Fourier Characterization of Sobolev Spaces}

Let $k$ be a nonnegative integer.

(i) A function $u \in L^2(\R^n)$ belongs to $H^k(\R^n)$ if and only if
$$(1 + |y|^k)\hat{u} \in L^2(\R^n).$$

(ii) In addition, there exists a positive constant $C$ such that
$$\frac{1}{C}\|u\|_{H^k(\R^n)} \le \|(1 + |y|^k)\hat{u}\|_{L^2(\R^n)} \le C\|u\|_{H^k(\R^n)}$$
for each $u \in H^k(\R^n)$.
\end{theorem}

\begin{proof}
\sketch For smooth compactly supported functions, Fourier transformation converts $D^\alpha$ into $(iy)^\alpha$. Plancherel and density give both norm bounds; conversely, inverse transforms of $(iy)^\alpha\hat u$ satisfy the weak derivative identities.
\end{proof}

\begin{definition}[Fractional Sobolev space $H^s(\R^n)$]
\label{def:fractional-sobolev-space}
\uses{thm:characterization-hk-fourier-transform}
\dcref{ch5:5.8-def-3}
\group{Fourier Characterization of Sobolev Spaces}

Assume $0 < s < \infty$ and $u \in L^2(\R^n)$. Then $u \in H^s(\R^n)$ if $(1 + |y|^s)\hat u \in L^2(\R^n)$. For noninteger $s$, we set
$$\|u\|_{H^s(\R^n)} := \|(1 + |y|^s)\hat u\|_{L^2(\R^n)}.$$
\end{definition}

\section{Other Spaces of Functions}
\label{sec:other-spaces-of-functions}

\subsection{The space $H^{-1}$}

\begin{definition}[The dual space $H^{-1}(U)$]
\label{def:dual-space-h-minus-1}
\uses{def:sobolev-space-zero-boundary}
\dcref{ch5:5.9-def-1}
\group{Dual Sobolev Spaces}

We denote by $H^{-1}(U)$ the dual space to $H_0^1(U)$.
\end{definition}

\begin{definition}[Norm on $H^{-1}(U)$]
\label{def:norm-on-h-minus-1}
\uses{def:dual-space-h-minus-1}
\dcref{ch5:5.9-def-2}
\group{Dual Sobolev Spaces}

If $f \in H^{-1}(U)$, we define the norm
$$\|f\|_{H^{-1}(U)} = \sup\left\{\langle f,u\rangle \mid u \in H_0^1(U), \|u\|_{H_0^1(U)} \le 1\right\}.$$
\end{definition}

\begin{theorem}[Characterization of $H^{-1}$]
\label{thm:characterization-h-minus-1}
\uses{def:norm-on-h-minus-1}
\dcref{ch5:5.9-thm-1}
\group{Dual Sobolev Spaces}

(i) Assume $f \in H^{-1}(U)$. Then there exist functions $f^0, f^1, \ldots, f^n$ in $L^2(U)$ such that
$$(1) \quad \langle f,v \rangle = \int_U f^0 v + \sum_{i=1}^n f^i v_{x_i}\,dx \quad (v \in H_0^1(U)).$$

(ii) Furthermore,
$$\|f\|_{H^{-1}(U)} = \inf\left\{\left(\int_U \sum_{i=0}^n |f^i|^2\,dx\right)^{1/2} \;\middle|\; f \text{ satisfies (1) for } f^0,\ldots,f^n \in L^2(U)\right\}.$$
\end{theorem}

\begin{proof}
\sketch Use Riesz representation in $H_0^1$ with the inner product $(u,v)=\int(Du\cdot Dv+uv)$. This yields an $L^2$ representation; Cauchy--Schwarz and testing with the representing function prove the minimal norm formula.
\end{proof}

\subsection{Spaces involving time}

\begin{definition}[$L^p(0,T;X)$]
\label{def:lp-time-banach-space}
\dcref{ch5:5.9-def-3}
\group{Bochner Spaces}
The space
$$L^p(0, T; X)$$
consists of all measurable functions $\mathbf{u} : [0, T] \to X$ with
\begin{itemize}
\item[(i)] $\|\mathbf{u}\|_{L^p(0,T;X)} := \left(\int_0^T \|\mathbf{u}(t)\|^p\,dt\right)^{1/p} < \infty$
\end{itemize}
for $1 \le p < \infty$, and
\begin{itemize}
\item[(ii)] $\|\mathbf{u}\|_{L^\infty(0,T;X)} := \text{ess sup}_{0 \le t \le T} \|\mathbf{u}(t)\| < \infty$.
\end{itemize}
\end{definition}

\begin{definition}[$C([0,T];X)$]
\label{def:continuous-time-banach-space}
\dcref{ch5:5.9-def-4}
\group{Bochner Spaces}
The space
$$C([0, T]; X)$$
comprises all continuous functions $\mathbf{u} : [0, T] \to X$ with
$$\|\mathbf{u}\|_{C([0,T];X)} := \max_{0 \le t \le T} \|\mathbf{u}(t)\| < \infty.$$
\end{definition}

\begin{definition}[Weak derivative for Banach-space-valued functions]
\label{def:weak-derivative-banach-valued}
\dcref{ch5:5.9-def-5}
\group{Bochner Spaces}
Let $\mathbf{u} \in L^1(0, T; X)$. We say $\mathbf{v} \in L^1(0, T; X)$ is the weak derivative of $\mathbf{u}$, written
$$\mathbf{u}' = \mathbf{v},$$
provided
$$\int_0^T \phi'(t)\mathbf{u}(t)\,dt = -\int_0^T \phi(t)\mathbf{v}(t)\,dt$$
for all scalar test functions $\phi \in C_c^\infty(0, T)$.
\end{definition}

\begin{definition}[$W^{1,p}(0,T;X)$ and $H^1(0,T;X)$]
\label{def:sobolev-space-time-banach}
\uses{def:lp-time-banach-space, def:weak-derivative-banach-valued}
\dcref{ch5:5.9-def-6}
\group{Bochner Spaces}

(i) The Sobolev space
$$W^{1,p}(0,T;X)$$
consists of all functions $u \in L^p(0,T;X)$ such that $u'$ exists in the weak sense and belongs to $L^p(0,T;X)$. Furthermore,
$$\|u\|_{W^{1,p}(0,T;X)} := \begin{cases}
\left(\int_0^T \|u(t)\|^p + \|u'(t)\|^p \, dt\right)^{1/p} & (1 \le p < \infty) \\
\text{ess sup}_{0 \le t \le T}(\|u(t)\| + \|u'(t)\|) & (p = \infty).
\end{cases}$$

(ii) We write $H^1(0,T;X) = W^{1,2}(0,T;X)$.
\end{definition}

\begin{theorem}[Calculus in an abstract space]
\label{thm:calculus-abstract-space}
\uses{def:sobolev-space-time-banach, thm:local-smooth-approximation}
\dcref{ch5:5.9-thm-2}
\group{Bochner Spaces}

Let $u \in W^{1,p}(0,T;X)$ for some $1 \le p \le \infty$. Then

(i) $u \in C([0,T];X)$ (after possibly being redefined on a set of measure zero), and

(ii) $u(t) = u(s) + \int_s^t u'(\tau) \, d\tau$ for all $0 \le s \le t \le T$.

(iii) Furthermore, we have the estimate
$$(7) \quad \max_{0 \le t \le T} \|u(t)\| \le C\|u\|_{W^{1,p}(0,T;X)},$$
the constant $C$ depending only on $T$.
\end{theorem}

\begin{proof}
\sketch Extend by zero, mollify in time, and pass to the limit in the classical fundamental theorem for the smooth convolutions. The integral representative is continuous and gives the stated norm estimate.
\end{proof}

\begin{theorem}[More calculus]
\label{thm:more-calculus-h1-hminus1}
\uses{def:sobolev-space-time-banach, def:dual-space-h-minus-1}
\dcref{ch5:5.9-thm-3}
\group{Bochner Spaces}

Suppose $u \in L^2(0,T;H_0^1(U))$, with $u' \in L^2(0,T;H^{-1}(U))$.

(i) Then
$$u \in C([0,T];L^2(U))$$
(after possibly being redefined on a set of measure zero).

(ii) The mapping
$$t \mapsto \|u(t)\|^2_{L^2(U)}$$
is absolutely continuous, with
$$\frac{d}{dt}\|u(t)\|^2_{L^2(U)} = 2\langle u'(t), u(t)\rangle$$
for a.e. $0 \le t \le T$.

(iii) Furthermore, we have the estimate
$$(10) \quad \max_{0 \le t \le T} \|u(t)\|_{L^2(U)} \le C\left(\|u\|_{L^2(0,T;H_0^1(U))} + \|u'\|_{L^2(0,T;H^{-1}(U))}\right),$$
the constant $C$ depending only on $T$.
\end{theorem}

\begin{proof}
\sketch Time mollification and the identity for the $L^2$ energy show the regularizations are Cauchy in $C([0,T];L^2)$. Passing to the limit gives absolute continuity of the energy and the duality formula; Cauchy--Schwarz yields the estimate.
\end{proof}

\begin{theorem}[Mappings into better spaces]
\label{thm:mappings-into-better-spaces}
\uses{def:sobolev-space-time-banach, thm:sobolev-extension-theorem, thm:difference-quotients-weak-derivatives, thm:more-calculus-h1-hminus1}
\dcref{ch5:5.9-thm-4}
\group{Bochner Spaces}

Assume that $U$ is open, bounded, and $\partial U$ is smooth. Take $m$ to be a nonnegative integer.

Suppose $\mathbf{u} \in L^2(0,T;H^{m+2}(U))$, with $\mathbf{u}' \in L^2(0,T;H^m(U))$.

(i) Then
$$\mathbf{u} \in C([0,T];H^{m+1}(U))$$
(after possibly being redefined on a set of measure zero).

(ii) Furthermore, we have the estimate
$$(13) \quad \max_{0 \le t \le T} \|\mathbf{u}(t)\|_{H^{m+1}(U)} \le C\bigl(\|\mathbf{u}\|_{L^2(0,T;H^{m+2}(U))} + \|\mathbf{u}'\|_{L^2(0,T;H^m(U))}\bigr),$$
the constant $C$ depending only on $T$, $U$, and $m$.
\end{theorem}

\begin{proof}
\sketch Extend in space and mollify in time. The $m=0$ energy estimate controls the $H^1$ norm uniformly; apply it to each spatial derivative of order at most $m$ and sum.
\end{proof}

\section{References}
\label{sec:sobolev-references}
\textbf{Sobolev-space references.} See \cite{GilbargTrudinger1983,LiebLoss1997,Ziemer1989,EvansGariepy1992,Temam1977}.
